% Zumdahl Chapter 9 test -- deliberately scaled down: only 9.1 is on the CED,
% so this is 15 MCQ + 1 FRQ (25 pts, 35 min) rather than the usual 20 + 2.
\documentclass{apchem-exam}

\apcunitword{Chapter}
\apcunit{9}
\apcunittitle{Covalent Bonding: Orbitals}
\apcdoctitle{Chapter Test}
\apcsubtitle{Zumdahl \S9.1 only \textbullet\ 35 minutes \textbullet\ shorter by design --- see scope note}

\begin{document}
\apctitleblock

\begin{apcnote}[Scope --- why this test is shorter]
Only \S9.1 of this chapter is on the AP framework, so this test is
\textbf{15 multiple choice $+$ 1 free response} rather than the usual
20 $+$ 2. Nothing from \S9.2--9.5 (molecular orbital theory, MO diagrams,
bond order from filled orbitals) appears --- that material was marked
enrichment in the notes and is not assessed.
\end{apcnote}

\examsection{I}{Multiple Choice}{15 questions \textbullet\ 22 minutes \textbullet\ 1 point each}{\nocalculator}

\begin{mcq}
  Hybrid orbitals are proposed in order to explain
  \choices
    \choice{why atoms have discrete energy levels}
    \correct{observed molecular geometries that atomic orbitals cannot
      account for}
    \choice{the existence of isotopes}
    \choice{why some molecules are coloured}
\end{mcq}
\why{Carbon's $2s^2 2p^2$ configuration predicts two bonds at 90$^\circ$;
methane has four at 109.5$^\circ$.}

\begin{mcq}
  The hybridization of the central atom is determined by the number of
  \choices
    \choice{bonds it forms}
    \correct{electron domains around it}
    \choice{valence electrons it has}
    \choice{atoms in the molecule}
\end{mcq}
\why{Lone pairs count as domains too --- which is why water is sp$^3$
despite forming only two bonds.}
\distractor{A}{the single most common error on this topic}

\begin{mcq}
  A central atom with three electron domains is
  \choices
    \choice{sp hybridized}
    \correct{sp$^2$ hybridized}
    \choice{sp$^3$ hybridized}
    \choice{not hybridized}
\end{mcq}
\why{Three domains require three hybrid orbitals: one $s$ plus two $p$.}

\begin{mcq}
  The hybridization of oxygen in \ce{H2O} is
  \choices
    \choice{sp}
    \choice{sp$^2$}
    \correct{sp$^3$}
    \choice{unhybridized}
\end{mcq}
\why{Two bonding pairs plus two lone pairs is four domains.}
\distractor{A}{counts only the two bonds}

\begin{mcq}
  How many hybrid orbitals result from sp$^2$ hybridization?
  \choices
    \choice{2}
    \correct{3}
    \choice{4}
    \choice{1}
\end{mcq}
\why{The number of hybrids always equals the number of atomic orbitals
combined --- here one $s$ and two $p$.}

\begin{mcq}
  An sp$^2$ hybridized carbon atom has how many unhybridized $p$ orbitals
  remaining?
  \choices
    \choice{0}
    \correct{1}
    \choice{2}
    \choice{3}
\end{mcq}
\why{Carbon has three $p$ orbitals; two are used in the hybrid.}

\begin{mcq}
  The hybridization of carbon in \ce{HCN} is
  \choices
    \correct{sp}
    \choice{sp$^2$}
    \choice{sp$^3$}
    \choice{dsp$^3$}
\end{mcq}
\why{Two domains --- one single bond and one triple bond.}

\begin{mcq}
  A $\sigma$ bond is formed by
  \choices
    \correct{head-on overlap along the internuclear axis}
    \choice{side-by-side overlap of $p$ orbitals}
    \choice{transfer of electrons}
    \choice{overlap of two lone pairs}
\end{mcq}
\why{That is the definition; $\pi$ bonds are the side-by-side case.}

\begin{mcq}
  A triple bond consists of
  \choices
    \choice{three $\sigma$ bonds}
    \choice{three $\pi$ bonds}
    \correct{one $\sigma$ and two $\pi$ bonds}
    \choice{two $\sigma$ and one $\pi$ bond}
\end{mcq}
\why{The first bond between two atoms is always $\sigma$; additional bonds
are $\pi$.}

\begin{mcq}
  How many $\sigma$ and $\pi$ bonds are in \ce{C2H4}?
  \choices
    \choice{4 $\sigma$, 2 $\pi$}
    \correct{5 $\sigma$, 1 $\pi$}
    \choice{6 $\sigma$, 0 $\pi$}
    \choice{5 $\sigma$, 2 $\pi$}
\end{mcq}
\why{Four C--H bonds plus the C=C $\sigma$, with one $\pi$ from the double
bond.}

\begin{mcq}
  Rotation about a carbon--carbon double bond is restricted because
  \choices
    \choice{the $\sigma$ bond is too strong}
    \correct{rotating would destroy the parallel overlap of the $\pi$ bond}
    \choice{the carbons are sp$^3$ hybridized}
    \choice{double bonds are longer than single bonds}
\end{mcq}
\why{The $p$ orbitals must stay parallel to overlap at all.}
\distractor{D}{double bonds are shorter, not longer}

\begin{mcq}
  \emph{Cis} and \emph{trans} isomers of a disubstituted alkene are
  separable compounds because
  \choices
    \correct{the $\pi$ bond prevents interconversion at ordinary
      temperatures}
    \choice{they contain different numbers of atoms}
    \choice{they have different molecular formulas}
    \choice{one is ionic}
\end{mcq}
\why{With a single bond the arrangements interconvert freely and are the
same substance.}

\begin{mcq}
  In \ce{CH3OH}, the hybridizations of carbon and oxygen are respectively
  \choices
    \correct{sp$^3$ and sp$^3$}
    \choice{sp$^3$ and sp$^2$}
    \choice{sp$^2$ and sp$^3$}
    \choice{sp and sp$^3$}
\end{mcq}
\why{Carbon has four bonds; oxygen has two bonds plus two lone pairs ---
four domains each.}

\begin{mcq}
  Which molecule contains an sp hybridized central atom?
  \choices
    \choice{\ce{NH3}}
    \choice{\ce{SO3}}
    \correct{\ce{CO2}}
    \choice{\ce{CH4}}
\end{mcq}
\why{Two domains gives sp; the others have three or four.}

\begin{mcq}
  \ce{SO2} and \ce{SO3} both have sp$^2$ hybridized sulfur, yet only
  \ce{SO2} is bent. The reason is that \ce{SO2}
  \choices
    \correct{has a lone pair occupying one of the three domains}
    \choice{has fewer electron domains}
    \choice{contains a different hybridization}
    \choice{has stronger bonds}
\end{mcq}
\why{Molecular geometry names only atom positions; the lone pair takes one
of the three trigonal planar positions.}

\examsection{II}{Free Response}{1 question \textbullet\ 13 minutes}{\calculatorok}

\frq{long}{10}{Zumdahl \S9.1 \textbullet\ hybridization and bonding}

Consider the molecule \ce{CH3-CH=CH-C#CH}.

\begin{parts}
  \item Give the hybridization of each of the five carbon atoms, reading
        left to right, with the domain count that justifies each.
        \answerlines[3]{C1 (\ce{CH3}): 4 domains $\to$ sp$^3$.
        C2 and C3 (the C=C pair): 3 domains each $\to$ sp$^2$.
        C4 and C5 (the C$\equiv$C pair): 2 domains each $\to$ sp.}
  \item State the total number of $\sigma$ bonds and $\pi$ bonds in the
        molecule. \answerlines[2]{Formula is \ce{C5H6}: 10 $\sigma$ bonds
        and 3 $\pi$ bonds.}
  \item Explain why the C1--C2 bond permits free rotation while the C2=C3
        bond does not. \answerlines[3]{C1--C2 is a single $\sigma$ bond,
        cylindrically symmetric about the axis, so rotation does not change
        the overlap. C2=C3 contains a $\pi$ bond whose two $p$ orbitals must
        remain parallel; rotating would set them perpendicular and destroy
        the overlap --- that is, break the bond.}
  \item Predict the C1--C2--C3 bond angle and the C3--C4--C5 bond angle.
        \answerlines[2]{About 120$^\circ$ at C2 (sp$^2$, trigonal planar)
        and 180$^\circ$ at C4 (sp, linear).}
\end{parts}

\begin{rubric}
  \rpoint{3}{(a) 1 pt sp$^3$ carbon, 1 pt both sp$^2$ carbons, 1 pt both sp
    carbons --- each requires the domain count; bare labels earn half}
  \rpoint{2}{(b) 1 pt $\sigma$ count, 1 pt $\pi$ count}
  \rpoint{3}{(c) 1 pt $\sigma$ cylindrical symmetry, 1 pt $\pi$ requires
    parallel $p$ orbitals, 1 pt rotating breaks the bond --- ``double bonds
    are stronger'' earns 0}
  \rpoint{2}{(d) 1 pt per angle, each consistent with the hybridization
    given in (a)}
\end{rubric}

\examtotal

\end{document}
